\chapter{Drift and gambler's ruin}
\label{ch:drift}

\begin{question}
A gambler wins each bet with probability $p$. What is the chance of reaching a target before going broke, and how long does it take? And what is the continuous version: Brownian motion with a drift?
\end{question}

\begin{whatfor}
\begin{itemize}
  \item Defines Brownian motion with drift $\mu$ and builds it from a slightly unfair random walk.
  \item Solves the gambler's ruin problem: win probability and expected duration.
  \item Finds the probability that Brownian motion with drift hits $+A$ before $-B$, and the chance it ever reaches $A$.
\end{itemize}
\skipline{fs:drift}
\end{whatfor}

\begin{prereq}
From Part~0: total probability, independence, geometric sums. From \crefrange{ch:first-mg}{ch:bm-walks}: Wiener process, exponential martingale $e^{uW_t-u^2t/2}$, random walk, optional stopping theorem and its exit recipe, the scaled random walk ($\Delta x=\sqrt{\Delta t}$).
\end{prereq}

%=====================================================================
\section{Brownian motion with drift}

\begin{definition}[Brownian motion with drift]
\label{def:drift}
A process $(B(t))_{t\ge0}$ is a \dfn{Brownian motion with drift} $\mu$ if (1) $B(0)=0$; (2) its paths are continuous a.s.; (3) it has stationary and independent increments; (4) $B(t)\sim N(\mu t,\,t)$. Equivalently $B(t)=\mu t+W(t)$ with $W$ a standard Brownian motion.
\end{definition}

It is a fair Brownian motion riding on a straight line of slope $\mu$.

\begin{pic}
\begin{tikzpicture}
\begin{axis}[width=10cm,height=4.6cm,xmin=0,xmax=1,xlabel={$t$},ylabel={$B(t)$},legend style={font=\scriptsize,at={(0.02,0.98)},anchor=north west}]
  \addplot[defc,thick] table[x=t,y expr=1.5*\thisrow{t}+\thisrow{w2}] {figs/bm3.dat}; \addlegendentry{$B(t)=1.5t+W(t)$}
  \addplot[prec,dashed,domain=0:1] {1.5*x}; \addlegendentry{mean $\mu t$}
\end{axis}
\end{tikzpicture}
\picnote{A Brownian motion with drift $\mu=1.5$ wanders around the line $\mu t$; its spread around the line grows like $\sqrt t$.}
\end{pic}

\paragraph{From an unfair random walk.} Take the scaled walk of \cref{ch:bm-walks} with $\Delta x=\sqrt{\Delta t}$, but let each step go right with probability $p=\frac12\big(1+\mu\sqrt{\Delta t}\big)$. Then each step $X_i=\pm1$ has
\[
  \E_P(X_i)=2p-1=\mu\sqrt{\Delta t},\qquad \Var(X_i)=1-(2p-1)^2=1-\mu^2\Delta t,
\]
so for $B(t)=\Delta x\,(X_1+\dots+X_{[t/\Delta t]})$,
\[
  \E_P\,B(t)=\sqrt{\Delta t}\,\Big[\frac{t}{\Delta t}\Big]\,\mu\sqrt{\Delta t}\to\mu t,\qquad
  \Var\,B(t)=\Delta t\,\Big[\frac{t}{\Delta t}\Big]\,(1-\mu^2\Delta t)\to t .
\]
A tiny bias in each coin, of order $\sqrt{\Delta t}$, adds up to a steady drift $\mu t$.

%=====================================================================
\section{Gambler's ruin}

\begin{setup}[The gambler]
A gambler has Rs~$i$. Each bet wins Rs~1 with probability $p$ and loses Rs~1 with probability $q=1-p$, independently. The gambler stops on reaching Rs~$N$ (the target) or Rs~$0$ (ruin). Write $X_n$ for the fortune after $n$ bets, $r=q/p$, and
\[
  P_i=P(\text{reach }N\text{ before }0\given X_0=i).
\]
\end{setup}

\begin{pic}
\begin{tikzpicture}[x=1.25cm,every node/.style={font=\small}]
  \foreach \i/\l in {0/0,1/1,2/2,4/N\!-\!2,5/N\!-\!1,6/N} \node[draw,circle,minimum size=1.05cm,inner sep=0pt,font=\scriptsize] (n\i) at (\i*1.2,0) {$\l$};
  \node at (3.6,0) {$\cdots$};
  \foreach \a/\b in {1/2,4/5} {\draw[-{Latex},bend left=35] (n\a) to node[above,font=\scriptsize]{$p$} (n\b); \draw[-{Latex},bend left=35] (n\b) to node[below,font=\scriptsize]{$q$} (n\a);}
  \draw[-{Latex},bend left=35] (n1) to node[below,font=\scriptsize]{$q$} (n0);
  \draw[-{Latex},bend left=35] (n5) to node[above,font=\scriptsize]{$p$} (n6);
  \node[trapc,below=0.45cm of n0] {ruin}; \node[thmc,below=0.45cm of n6] {target};
\end{tikzpicture}
\picnote{The fortune moves one step right with probability $p$, left with probability $q$. The states $0$ and $N$ are absorbing: once there, the game is over.}
\end{pic}

\begin{theorem}[Gambler's ruin]
\label{thm:ruin}
For $0\le i\le N$,
\[
  P_i=\begin{cases}\dfrac{1-r^i}{1-r^N},&p\ne\frac12,\\[2ex] \dfrac iN,&p=\frac12.\end{cases}
\]
\end{theorem}
\begin{proof}
Condition on the first bet (first-step analysis).
\begin{steps}
  \item \emph{Boundaries:} $P_0=0$, $P_N=1$.
  \item \emph{First bet:} after it the fortune is $i+1$ or $i-1$, and the rest is a fresh game from there. By total probability, $P_i=p\,P_{i+1}+q\,P_{i-1}$ for $1\le i\le N-1$.
  \item \emph{Differences:} write $P_i=(p+q)P_i$ and rearrange: $p(P_{i+1}-P_i)=q(P_i-P_{i-1})$, so $P_{i+1}-P_i=r\,(P_i-P_{i-1})$. Repeating, and using $P_1-P_0=P_1$,
  \[
    P_{k+1}-P_k=r^k\,P_1 .
  \]
  \item \emph{Telescope:} adding the differences for $k=0,\dots,i-1$,
  \[
    P_i=(1+r+\dots+r^{i-1})P_1=\begin{cases}\frac{1-r^i}{1-r}P_1,&r\ne1,\\ i\,P_1,&r=1.\end{cases}
  \]
  \item \emph{Use $P_N=1$:} $P_1=\frac{1-r}{1-r^N}$ (or $\frac1N$ when $r=1$, i.e.\ $p=\frac12$). Substitute. \qedhere
\end{steps}
\end{proof}

\begin{result}[Infinite bankroll of the opponent]
\label{res:ruin-limit}
Let $N\to\infty$ (the gambler plays until ruin, aiming at nothing). Then
\[
  P(\text{gambler is never ruined})=\lim_{N\to\infty}P_i=\begin{cases}0,&p\le\frac12,\\ 1-\big(\frac qp\big)^i,&p>\frac12.\end{cases}
\]
\end{result}
\begin{proof}
If $p<\frac12$ then $r>1$ and $r^N\to\infty$, so $P_i\to0$. If $p=\frac12$, $\frac iN\to0$. If $p>\frac12$ then $r<1$, $r^N\to0$, and $P_i\to1-r^i$.
\end{proof}

A fair game against an infinitely rich opponent ruins you for sure.

\begin{result}[Expected duration]
\label{res:duration}
Let $D_i=\E_P(\tau\given X_0=i)$, where $\tau$ is the number of bets played. Then
\[
  D_i=\begin{cases}\dfrac{N\,P_i-i}{p-q},&p\ne\frac12,\\[2ex] i(N-i),&p=\frac12.\end{cases}
\]
\end{result}
\begin{proof}
Use optional stopping with two martingales; the stopped walk stays in $[0,N]$.
\begin{steps}
  \item Each bet has mean $p-q$, so $X_n-n(p-q)$ is a martingale (\cref{ch:martingales}). At $\tau$, $X_\tau=N$ with probability $P_i$, else $0$. Optional stopping: $N P_i-(p-q)D_i=i$; solve for $D_i$.
  \item For $p=\frac12$ that equation is empty ($0\cdot D_i$); use the martingale $(X_n-i)^2-n$ instead, as in the exit recipe of \cref{ch:stopping}: $D_i=\E_P(X_\tau-i)^2=(N-i)^2\frac iN+i^2\frac{N-i}{N}=i(N-i)$. \qedhere
\end{steps}
\end{proof}

\begin{example}[Roulette]
\label{ex:roulette}
Betting on red with 38 pockets: $p=\frac{18}{38}\approx0.474$, $q\approx0.526$, so $r=\frac{20}{18}=\frac{10}{9}$. Start with $i=10$, target $N=20$:
\[
  P_{10}=\frac{1-(10/9)^{10}}{1-(10/9)^{20}}=\frac{1}{1+(10/9)^{10}}\approx0.259,\qquad
  D_{10}=\frac{20(0.259)-10}{18/38-20/38}\approx91.8 .
\]
A fair coin would give $P_{10}=\frac12$ and $D_{10}=100$. A small house edge cuts the chance of doubling your money almost in half.
\end{example}

%=====================================================================
\section{Exit probabilities for Brownian motion with drift}

\begin{theorem}[Up $A$ before down $B$]
\label{thm:exit-drift}
Let $B(t)=\mu t+W(t)$ with $\mu\ne0$, and $A,B>0$. Then
\[
  P\big(B\text{ hits }A\text{ before }-B\big)=\frac{1-e^{-2\mu B}}{1-e^{-2\mu(A+B)}}=\frac{e^{2\mu B}-1}{e^{2\mu B}-e^{-2\mu A}} .
\]
For $\mu=0$ it is $\frac{B}{A+B}$.
\end{theorem}
\begin{proof}
Find a martingale that is a function of $B(t)$ alone.
\begin{steps}
  \item $e^{-2\mu B(t)}=e^{-2\mu W(t)-2\mu^2t}=e^{uW(t)-\frac12u^2t}$ with $u=-2\mu$: the exponential martingale of \cref{ch:first-mg}.
  \item Let $\tau$ be the exit time and $p$ the probability of exiting at $A$. The stopped martingale is bounded, so optional stopping gives $1=\E_Pe^{-2\mu B(\tau)}=p\,e^{-2\mu A}+(1-p)\,e^{2\mu B}$.
  \item Solve: $p=\dfrac{e^{2\mu B}-1}{e^{2\mu B}-e^{-2\mu A}}$; multiply top and bottom by $e^{-2\mu B}$ for the first form.
  \item For $\mu=0$, $B(t)=W(t)$ and the recipe of \cref{ch:stopping} gives $\frac{B}{A+B}$ (also the limit $\mu\to0$). \qedhere
\end{steps}
\end{proof}

\begin{remark}[Second route: the random-walk limit]
The same answer comes from the gambler's ruin. Put the walk of the first section on the grid $-B,\dots,A$ with step $\Delta x$: start $i=B/\Delta x$, target $N=(A+B)/\Delta x$, and $r=\frac qp=\frac{1-\mu\Delta x}{1+\mu\Delta x}$. Taking logarithms and using $\ln(1\pm h)=\pm h+O(h^2)$,
\[
  \ln r^{1/\Delta x}=\frac{\ln(1-\mu\Delta x)-\ln(1+\mu\Delta x)}{\Delta x}\to-2\mu,\qquad\text{so } r^{B/\Delta x}\to e^{-2\mu B}.
\]
Then $P_i=\frac{1-r^{B/\Delta x}}{1-r^{(A+B)/\Delta x}}\to\frac{1-e^{-2\mu B}}{1-e^{-2\mu(A+B)}}$.
\end{remark}

\begin{result}[Will it ever reach $A$?]
\label{res:ever}
For $\mu<0$ and $A>0$, $P\big(B(t)\text{ ever reaches }A\big)=e^{2\mu A}$. Equivalently, the overall maximum $\max_{t\ge0}B(t)$ is exponential with rate $-2\mu$. For $\mu\ge0$ the probability is $1$.
\end{result}
\begin{proof}
Let $B\to\infty$ in \cref{thm:exit-drift}. For $\mu<0$, $e^{-2\mu B}\to\infty$, and $\frac{1-e^{-2\mu B}}{1-e^{-2\mu(A+B)}}\to\frac{e^{-2\mu B}}{e^{-2\mu(A+B)}}=e^{2\mu A}$. For $\mu>0$ the ratio tends to $1$; for $\mu=0$, $\frac{B}{A+B}\to1$.
\end{proof}

\begin{pic}
\begin{tikzpicture}
\begin{axis}[width=9cm,height=4.6cm,xmin=-2,xmax=2,ymin=0,ymax=1,xlabel={drift $\mu$},ylabel={$P(\text{hit }1\text{ before }-1)$},samples=101,domain=-2:2]
  \addplot[defc,very thick] {(abs(x)<1e-6) ? 0.5 : (1-exp(-2*x))/(1-exp(-4*x))};
  \addplot[prec,dashed] coordinates {(-2,0.5)(2,0.5)};
\end{axis}
\end{tikzpicture}
\picnote{With $A=B=1$ the exit probability is $\frac{1}{1+e^{-2\mu}}$: $\frac12$ for no drift, and quickly close to $0$ or $1$ as the drift grows.}
\end{pic}

\begin{example}[When to sell?]
\label{ex:sell}
A stock price, measured from today's level, moves as $B(t)=\mu t+W(t)$ with $\mu=-d<0$. You hold the right to buy at price $A>0$ (above today's level) and will exercise when the price first reaches some level $x>A$, gaining $x-A$. Which $x$ maximises the expected gain?
\begin{steps}
  \item $P(\text{price ever reaches }x)=e^{2\mu x}=e^{-2dx}$, so the expected gain is $f(x)=(x-A)e^{-2dx}$.
  \item $f'(x)=e^{-2dx}\big[1-2d(x-A)\big]=0$ gives $x^*=A+\frac{1}{2d}$.
  \item Maximum expected gain: $f(x^*)=\frac{1}{2d}e^{-2dA-1}$. With $d=0.5$, $A=1$: $x^*=2$ and $f(x^*)=e^{-2}\approx0.135$.
\end{steps}
\end{example}

%=====================================================================
\formulasheet{fs:drift}

Exits, the gambler's ruin and ``will it ever reach $b$'' are rows 1--4 of Card~\ref{card:ost} (beneath \cref{ch:stopping}), with the martingales of \cref{card:incr}, row 6. This sheet holds the law of the process itself.

\begin{mastercard}[label=card:Ylaw]{The law of $Y_t=x_0+\mu t+\sigma W_t$ given any known values}
Only the nearest known value on each side of $t$ matters. With known values $Y_{s_1}=a_1$ before $t$ and $Y_{s_2}=a_2$ after $t$ ($s_1<t<s_2$):
\[
  \boxed{\;Y_t\given(Y_{s_1}=a_1,\,Y_{s_2}=a_2)\ \sim\ N\Big(\underbrace{\frac{s_2-t}{s_2-s_1}}_{\text{weight on }a_1}a_1+\underbrace{\frac{t-s_1}{s_2-s_1}}_{\text{weight on }a_2}a_2,\ \ \sigma^2\,\frac{(t-s_1)(s_2-t)}{s_2-s_1}\Big).\;}
\]
The mean is the straight line through the two known points; the variance is $\sigma^2\times$(distance to the left point)$\times$(distance to the right point)$\div$(gap), zero at the known points and largest in the middle. The drift $\mu$ does \emph{not} appear: it is already inside $a_1$ and $a_2$. When nothing is known after $t$, use the forward rows, where $\mu$ \emph{does} appear.
\begin{plugrows}
  \plugrow{Only the start is known}{\pstep{Box: mean $\frac{s_2-t}{s_2-s_1}a_1+\frac{t-s_1}{s_2-s_1}a_2$, variance $\sigma^2\frac{(t-s_1)(s_2-t)}{s_2-s_1}$.}\pstep{Nothing is known after $t$: use the forward law instead (row 2) from the start.}\pstep{From the start: $s=0$, $a=x_0$.}\pstep{Substitute into row 2: $N\big(x_0+\mu t,\ \sigma^2t\big)$.}}{$Y_t\sim N(x_0+\mu t,\ \sigma^2t)$}{always}{$Y=1.5t+W$ ($x_0=0$, $\mu=1.5$, $\sigma=1$), $t=2$: $Y_2\sim N(3,2)$; $P(Y_2<0)=\Phi\big(\frac{0-3}{\sqrt2}\big)=\Phi(-2.12)\approx0.017$}
  \plugrow{Forward: latest value $Y_s=a$ known, $t>s$}{\pstep{Split (\cref{card:freeze}): $Y_t=Y_s+(Y_t-Y_s)$; $Y_s=a$ known, $Y_t-Y_s\sim N(\mu(t-s),\sigma^2(t-s))$ fresh.}\pstep{Known plus fresh normal: shift the fresh normal by $a$.}\pstep{Earlier values are in the past and add nothing.}}{$N\big(a+\mu(t-s),\ \sigma^2(t-s)\big)$}{$t$ after every known time}{race $Y=\sigma W$, $Y_{1/2}=\sigma$, $t=1$: $Y_1\sim N(\sigma,\frac{\sigma^2}2)$; $P(Y_1>0)=\Phi\big(\frac{\sigma}{\sigma/\sqrt2}\big)=\Phi(\sqrt2)\approx0.921$}
  \plugrow{Backward: start $x_0$ and a later value $Y_{s_2}=a_2$}{\pstep{Box: mean $\frac{s_2-t}{s_2-s_1}a_1+\frac{t-s_1}{s_2-s_1}a_2$, variance $\sigma^2\frac{(t-s_1)(s_2-t)}{s_2-s_1}$.}\pstep{From the start: $s_1=0$, $a_1=x_0$.}\pstep{Substitute: weights $\frac{s_2-t}{s_2}$ on $x_0$ and $\frac t{s_2}$ on $a_2$; variance $\sigma^2\frac{t(s_2-t)}{s_2}$.}}{$N\Big(x_0+\frac{t}{s_2}(a_2-x_0),\ \sigma^2\frac{t(s_2-t)}{s_2}\Big)$}{$t<s_2$}{$Y=3W$, $Y_{10}=12$, $t=4$: mean $\frac4{10}\cdot12=4.8$, variance $9\cdot\frac{4\cdot6}{10}=21.6$; race, $Y_1=\sigma$, $t=\frac12$: $N(\frac\sigma2,\frac{\sigma^2}4)$, $P(Y_{1/2}>0)=\Phi\big(\frac{\sigma/2}{\sigma/2}\big)=\Phi(1)\approx0.841$}
  \plugrow{Between two known values}{\pstep{Box with the nearest known value on each side: $(s_1,a_1)$ before $t$, $(s_2,a_2)$ after.}\pstep{Weights: $\frac{s_2-t}{s_2-s_1}$ and $\frac{t-s_1}{s_2-s_1}$.}\pstep{Variance: $\sigma^2\times$(distance left)$\times$(distance right)$\div$(gap).}}{the box}{$s_1<t<s_2$}{$W_1=2$, $W_5=10$, $t=2$: weights $\frac{5-2}4=\frac34$, $\frac{2-1}4=\frac14$, mean $\frac34(2)+\frac14(10)=4$, variance $\frac{1\cdot3}4=0.75$; $Y=2t+3W$, $Y_1=5$, $Y_5=13$: mean $\frac34(5)+\frac14(13)=7$, variance $9(0.75)=6.75$}
  \plugrow{Pinned at both ends (bridge)}{\pstep{Box: mean $\frac{s_2-t}{s_2-s_1}a_1+\frac{t-s_1}{s_2-s_1}a_2$, variance $\sigma^2\frac{(t-s_1)(s_2-t)}{s_2-s_1}$.}\pstep{From the bridge: $s_1=0$, $a_1=0$, $s_2=1$, $a_2=0$.}\pstep{Substitute: mean $0$, variance $\sigma^2\frac{t(1-t)}1$.}}{$N\big(0,\ \sigma^2t(1-t)\big)$}{$0<t<1$}{$W$ given $W_1=0$, $t=\frac12$: variance $\frac12\cdot\frac12=\frac14$, so $N(0,\frac14)$}
  \plugrow{Several times at once (joint density)}{\pstep{Row 2 for each step: given $Y_{t_{k-1}}=y_{k-1}$, $Y_{t_k}\sim N\big(y_{k-1}+\mu(t_k-t_{k-1}),\sigma^2(t_k-t_{k-1})\big)$.}\pstep{\cref{card:bayes}, density form: joint density = product of these conditional densities.}}{$f(y_1,\dots,y_n)=\prod_kf_{\sigma^2(t_k-t_{k-1})}\big(y_k-y_{k-1}-\mu(t_k-t_{k-1})\big)$, $f_v(z)=\frac{e^{-z^2/2v}}{\sqrt{2\pi v}}$}{$t_0=0$, $y_0=x_0$, times increasing}{$(W_1,W_3)$: $f_1(x)\,f_{3-1}(y-x)=f_1(x)f_2(y-x)$}
  \plugrow{A probability}{\pstep{Take the normal law $N(m,v)$ from a row above.}\pstep{\cref{card:normal-prob}: $P(Y>c)=1-\Phi\big(\frac{c-m}{\sqrt v}\big)$.}}{standardise}{divide by the standard deviation}{$W_2$ given $W_1=2$, $W_5=10$: $N(4,0.75)$ (row 4); $P(W_2>5)=1-\Phi\big(\frac{5-4}{\sqrt{0.75}}\big)=1-\Phi(1.155)\approx0.124$}
\end{plugrows}
\end{mastercard}
